\documentclass[10pt,a4paper]{amsart}
\usepackage[margin=19mm]{geometry}
\usepackage{amsmath,amssymb,mathtools}
\usepackage[scr=rsfs,cal=euler]{mathalfa}
\usepackage{xfrac}
\usepackage[unicode,hidelinks,bookmarks=false]{hyperref}
\newcommand{\ie}{i.e.\ }
\newcommand{\cf}{cf.\ }
\newcommand{\resp}{respectively\ }
\newcommand{\Subsection}{\subsection}
\providecommand{\nobreakdash}{\nobreak\hbox{-}}
\setlength{\emergencystretch}{2em}
\multlinegap0pt
\newcommand{\G}{\mathrm{Gr}}
\usepackage{amssymb}
\usepackage{mathtools}
\usepackage[shortlabels]{enumitem}
\usepackage{xcolor}
\newtheorem{lemma}{Lemma}
\newcommand{\PGL}{\mathrm{PGL}}
\title{On the Codimension-1 \texorpdfstring{$\PGL_4$}{PGL4} Orbit Closures in \texorpdfstring{$\G(2,10)$}{Gr(2,10)}}
\author{Ari Krishna}
\date{Source: arXiv:2603.28619v1 (CC BY 4.0)}
\begin{document}
\maketitle

\noindent\emph{Source and licence.} On the Codimension-1 \texorpdfstring{$\PGL_4$}{PGL4} Orbit Closures in \texorpdfstring{$\G(2,10)$}{Gr(2,10)}. Version arXiv:2603.28619v1 (CC BY 4.0), distributed by arXiv under the Creative Commons Attribution 4.0 International licence, http://creativecommons.org/licenses/by/4.0/, as recorded in the arXiv licence metadata for this exact version and shown on the abstract page https://arxiv.org/abs/2603.28619v1. Full licence notice: \texttt{licenses/arxiv-2603.28619-CC-BY-4.0.txt}.\par\smallskip

\noindent\textit{Editorial scope.} Counted unit: the article's lemma lem:normal-form (source0.tex lines 459--472). The article's complete original proof of it is delivered with it (source0.tex lines 474--508, one \texttt{proof} environment of 35 lines). No mathematics is written, repaired or re-derived: the statement and the proof are the article's own lines, in the article's own notation, and no retained line is substituted or re-flowed.  No further source-local context is needed: every cross-reference of every retained line resolves inside the two delivered spans.  Nothing was omitted from any retained span, so the package carries no omission record.  No retained line contains a citation, so the wrapper carries no bibliography and no bibliography entry.  No retained line contains a figure environment, an image-inclusion command or drawing code, so no image is delivered and the package declares no asset dependency.  Editorial formatting only, in addition: an AMS \texttt{amsart} preamble that restates the article's own declarations and macros used by the retained lines, namely the environments \texttt{lemma} on separate counters, together with the standard packages those lines need.  The retained environments print in this document as Lemma 1 in order of appearance, whereas the article prints the same items with the numbering of its own full document.  The complete native source of the article as distributed in the arXiv e-print for this exact version is bundled unchanged as \texttt{sources/197460/source0.tex}.
\begin{lemma}\label{lem:normal-form}
Let $W\in T$ be a general point. Then, $W$ is $\PGL(V)$-equivalent to a pencil of the form
\[
W_{a,b}:=
\left \langle
x_0^2+x_1^2+x_2^2,\
 x_0x_3 + a x_1^2 + b x_2^2
\right \rangle
\]
with $a,b\in k^\times$ and $a\neq b$. Its discriminant quartic has root type $2+1+1$, namely,
\[
\Delta_{W_{a,b}}(1,t) = -\frac14 t^2(1+at)(1+bt).
\]
\end{lemma}

\begin{proof}
Let $W=\langle Q_0,Q_1\rangle$ with $\ell_W$ tangent to $D$ at a smooth point $[Q_0]\in D$. Since $[Q_0]\in D_{\mathrm{sm}}$, the quadric $Q_0$ has rank $3$, and after a projective change of coordinates if necessary, we may assume that
\[
Q_0=x_0^2+x_1^2+x_2^2.
\]
Thus, $\ker(Q_0)=\langle e_3\rangle$.

Write
\[
Q_1=L(x_0,x_1,x_2)x_3 + q(x_0,x_1,x_2) + c x_3^2.
\]
The condition of tangency of the line $\ell_W$ to $D$ at $[Q_0]$ means that the linear term of
$\det(Q_0+tQ_1)$
in $t$ vanishes. Writing $A=\operatorname{diag}(1,1,1,0)$ for the matrix of $Q_0$, we have
\[
\frac{d}{dt}\det(A+tB)\big|_{t=0}=\operatorname{tr}(\operatorname{adj}(A)B),
\]
and here $\operatorname{adj}(A)$ is the matrix with a single $1$ in the $(4,4)$-entry. Therefore, the tangency condition translates exactly to the vanishing of the $x_3^2$-coefficient of $Q_1$, i.e. $c=0$. Moreover, $L\neq 0$, as otherwise every quadric in the pencil would be singular at $[0:0:0:1]$ and the line would lie in $D$.

By using the orthogonal group of the quadratic form $x_0^2+x_1^2+x_2^2$, we may arrange that $L=x_0$. Replacing $x_3$ by $x_3+\ell(x_0,x_1,x_2)$ leaves $Q_0$ unchanged, and it transforms
\[
Q_1=x_0x_3+q
\]
into
\[
x_0x_3+(x_0\ell+q).
\]
Choosing $\ell=-\alpha x_0-\beta x_1-\gamma x_2$ cancels the $x_0^2$, $x_0x_1$, and $x_0x_2$ terms of $q$. Finally, diagonalizing the remaining quadratic form in $x_1,x_2$ gives the desired normal form.

A direct determinant computation shows that
\[
\det(Q_0+tQ_1) = -\frac14 t^2(1+at)(1+bt).
\]
For a general tangent pencil, the two residual roots are distinct and nonzero, so $a,b\neq 0$ and $a\neq b$; the proposition follows.
\end{proof}

\end{document}
